\documentclass[11pt,a4paper]{article}

\usepackage[utf8]{inputenc}
\usepackage[english]{babel}
\usepackage{fontspec}
\setmainfont{Times New Roman}
\setsansfont{Arial}
\setmonofont{Consolas}

\usepackage[a4paper,top=1.8cm,bottom=2cm,left=1.8cm,right=1.8cm]{geometry}
\usepackage{amsmath,amssymb}
\usepackage{booktabs}
\usepackage{tabularx}
\usepackage{array}
\usepackage{listings}
\usepackage{xcolor}
\usepackage{fancyhdr}
\usepackage{multicol}
\usepackage{enumitem}
\usepackage{tcolorbox}
\usepackage{lastpage}

% Colors for syntax highlighting and frames
\definecolor{codegray}{rgb}{0.96,0.96,0.96}
\definecolor{keywordblue}{rgb}{0.0,0.2,0.6}
\definecolor{stringgreen}{rgb}{0.0,0.5,0.2}
\definecolor{commentgray}{rgb}{0.4,0.4,0.4}
\definecolor{correctgreen}{rgb}{0.0,0.5,0.15}
\definecolor{boxborder}{rgb}{0.7,0.7,0.7}

\lstset{
  language=Python,
  backgroundcolor=\color{codegray},
  basicstyle=\ttfamily\small,
  keywordstyle=\color{keywordblue}\bfseries,
  stringstyle=\color{stringgreen},
  commentstyle=\color{commentgray}\itshape,
  showstringspaces=false,
  breaklines=true,
  frame=single,
  rulecolor=\color{boxborder},
  aboveskip=4pt,
  belowskip=4pt,
  numbers=none,
  columns=fullflexible
}

% Header & Footer
\pagestyle{fancy}
\fancyhf{}
\setlength{\headheight}{15pt}
\lhead{\small \textbf{Course:} Data Analysis with Python (DSAI1005)}
\rhead{\small \textbf{SOLUTIONS \& GRADING RUBRIC (WEEKS 1 -- 4)}}
\lfoot{\small Faculty of Data Science \& AI --- National Economics University}
\rfoot{\small Page \thepage/\pageref{LastPage}}
\renewcommand{\headrulewidth}{0.5pt}
\renewcommand{\footrulewidth}{0.5pt}

\begin{document}

\noindent
\begin{minipage}[t]{0.55\textwidth}
  \centering
  \textbf{NATIONAL ECONOMICS UNIVERSITY}\\
  \textbf{FACULTY OF DATA SCIENCE \& AI}\\
  \rule{4.5cm}{0.4pt}\\[2mm]
  \textbf{OFFICIAL SOLUTIONS \& GRADING GUIDE}\\
  \textbf{Course:} Data Analysis with Python (DSAI1005)\\
  \textbf{Content:} Midterm Exam (Weeks 1 -- 4) --- \textbf{Exam Code:} \textbf{201-EN}
\end{minipage}
\hfill
\begin{minipage}[t]{0.42\textwidth}
  \begin{tcolorbox}[colback=blue!5!white,colframe=blue!75!black,title=\textbf{Score Breakdown (10-Point Scale)}]
    \small
    \textbf{Section I (Multiple-Choice):} 20 items $\times$ 0.25 pts = \textbf{5.0 pts}\\
    \textbf{Section II (Short-Answer):} 4 questions = \textbf{5.0 pts}\\
    \rule{\linewidth}{0.3pt}\\
    \textbf{Total Exam Score:} \textbf{10.0 points}
  \end{tcolorbox}
\end{minipage}

\vspace{3mm}
\hrule
\vspace{2mm}

\section*{SECTION I: MULTIPLE-CHOICE KEY (5.0 pts --- 0.25 pts each)}

\begin{center}
\begin{tabular}{|c|c|c|c|c|c|c|c|c|c|c|}
  \hline
  \textbf{Q} & \textbf{1} & \textbf{2} & \textbf{3} & \textbf{4} & \textbf{5} & \textbf{6} & \textbf{7} & \textbf{8} & \textbf{9} & \textbf{10} \\
  \hline
  \textbf{Key} & \textbf{B} & \textbf{C} & \textbf{A} & \textbf{B} & \textbf{B} & \textbf{B} & \textbf{A} & \textbf{B} & \textbf{B} & \textbf{B} \\
  \hline
  \hline
  \textbf{Q} & \textbf{11} & \textbf{12} & \textbf{13} & \textbf{14} & \textbf{15} & \textbf{16} & \textbf{17} & \textbf{18} & \textbf{19} & \textbf{20} \\
  \hline
  \textbf{Key} & \textbf{B} & \textbf{B} & \textbf{B} & \textbf{C} & \textbf{B} & \textbf{B} & \textbf{A} & \textbf{B} & \textbf{B} & \textbf{B} \\
  \hline
\end{tabular}
\end{center}

\subsection*{Concise Conceptual Explanations:}
\begin{itemize}[leftmargin=*,itemsep=1.5pt]
  \item \textbf{Q1 (B):} EDA (pioneered by John Tukey) seeks to explore data structure, test assumptions, discover patterns, and spot anomalies before modeling.
  \item \textbf{Q2 (C):} Scale property of standard deviation: multiplying all values by a constant $c = 3$ scales the standard deviation by $|c| = 3$ ($\text{Std}(cX) = |c|\text{Std}(X)$).
  \item \textbf{Q3 (A):} Tukey's standard box plot defines outer fences as $[Q_1 - 1.5\times\text{IQR}, \; Q_3 + 1.5\times\text{IQR}]$.
  \item \textbf{Q4 (B):} $r = 0$ indicates the absence of linear correlation; it does not preclude quadratic or other deterministic non-linear dependencies.
  \item \textbf{Q5 (B):} \texttt{np.arange(start=3, stop=15, step=4)} generates $[3, 7, 11]$ (stop value 15 is excluded).
  \item \textbf{Q6 (B):} \texttt{m[::-1, 1]} reverses the row order: $[7, 8, 9] \rightarrow [4, 5, 6] \rightarrow [1, 2, 3]$. Extracting column index 1 yields $[8, 5, 2]$.
  \item \textbf{Q7 (A):} Broadcasting rules: Dim 1: 5 vs 1 $\rightarrow 5$; Dim 2: 1 vs 4 $\rightarrow 4$; Dim 3: 3 vs 1 (implicit trailing) $\rightarrow 3$ $\Rightarrow$ shape $(5, 4, 3)$.
  \item \textbf{Q8 (B):} \texttt{np.where(condition, x, y)} is the vectorized equivalent of the ternary operator.
  \item \textbf{Q9 (B):} \texttt{axis=0} computes the sum down the rows (collapsing rows, preserving columns), producing a 1D array of shape $(5,)$.
  \item \textbf{Q10 (B):} A \texttt{pd.Series} is a 1D labeled array capable of holding any data type.
  \item \textbf{Q11 (B):} Label-based slicing with \texttt{.loc} is inclusive of both the start and stop boundaries $\rightarrow$ includes \texttt{'a'}, \texttt{'b'}, and \texttt{'c'}.
  \item \textbf{Q12 (B):} \texttt{df.fillna()} replaces missing entries (\texttt{NaN}) with designated constant values or summary statistics.
  \item \textbf{Q13 (B):} In Pandas, multiple boolean conditions require bitwise operators (\texttt{|}) and explicit parentheses around each clause.
  \item \textbf{Q14 (C):} The \texttt{'IT'} group has salaries 50 and 70; the maximum value is 70.
  \item \textbf{Q15 (B):} An inner join retains only those rows where the join key is mutually present in both DataFrames.
  \item \textbf{Q16 (B):} \texttt{plt.subplots(2, 3)} creates a $2 \times 3$ grid of 6 individual \texttt{Axes} subplots.
  \item \textbf{Q17 (A):} \texttt{sns.kdeplot()} estimates and visualizes the smooth probability density function of continuous variables.
  \item \textbf{Q18 (B):} \texttt{sns.pairplot()} renders pairwise scatter plots across numeric columns with histograms/KDEs along the diagonal.
  \item \textbf{Q19 (B):} \texttt{plt.tight\_layout()} automatically adjusts subplot padding to prevent overlapping titles, labels, and ticks.
  \item \textbf{Q20 (B):} \texttt{sns.barplot()} plots group means with bootstrapped confidence intervals (error bars) by default.
\end{itemize}

\vspace{3mm}
\hrule
\vspace{2mm}

\section*{SECTION II: SOLUTIONS \& GRADING RUBRIC (5.0 pts)}

\begin{tcolorbox}[colback=gray!5!white,colframe=gray!60!black,title=\textbf{General Grading Policy for Evaluators}]
\small
\begin{itemize}[leftmargin=*,itemsep=1pt]
  \item Do not penalize students for interchange of single quotes \texttt{' '} and double quotes \texttt{" "}.
  \item Minor trailing syntax slips (e.g., missing closing parenthesis at the line end) should not incur deductions if the algorithmic logic is sound.
  \item Mathematically and logically equivalent implementations in NumPy / Pandas / Matplotlib are entitled to full credit.
\end{itemize}
\end{tcolorbox}

\subsection*{Question 1: NumPy Vectorization \& Slicing (1.0 pt)}
\begin{enumerate}[label=\textbf{\alph*)}]
  \item \textbf{(0.5 pts):}
  \begin{itemize}
    \item \textbf{Mechanism:} Array **Broadcasting**. \textit{(0.25 pts)}
    \item \textbf{Resulting Matrix:} \textit{(0.25 pts)}
    $$\text{weighted\_scores} = \begin{bmatrix}
      8.0\times 0.2 & 7.0\times 0.3 & 9.0\times 0.5 \\
      6.0\times 0.2 & 8.0\times 0.3 & 7.0\times 0.5 \\
      9.0\times 0.2 & 9.0\times 0.3 & 8.5\times 0.5 \\
      5.0\times 0.2 & 6.0\times 0.3 & 6.5\times 0.5
    \end{bmatrix} = \begin{bmatrix}
      1.6 & 2.1 & 4.5 \\
      1.2 & 2.4 & 3.5 \\
      1.8 & 2.7 & 4.25 \\
      1.0 & 1.8 & 3.25
    \end{bmatrix}$$
  \end{itemize}

  \item \textbf{(0.5 pts):}
  \begin{itemize}
    \item \textbf{NumPy statement for gpa:} \textit{(0.25 pts)}
\begin{lstlisting}
gpa = np.sum(weighted_scores, axis=1)  # Or: weighted_scores.sum(axis=1)
\end{lstlisting}
    \item Row-wise sums yield: $[8.2, \; 7.1, \; 8.75, \; 6.05]$.
    \item Condition \texttt{gpa >= 8.0} is true for indices 0 and 2.
    \item Extracting column 0 (\texttt{scores[..., 0]}): \textbf{Printed output:} \texttt{[8.0, 9.0]} (or \texttt{[8., 9.]}). \textit{(0.25 pts)}
  \end{itemize}
\end{enumerate}

\subsection*{Question 2: Data Cleaning \& Aggregation with Pandas (1.5 pts)}
\begin{enumerate}[label=\textbf{\alph*)}]
  \item \textbf{(0.5 pts):} Filtering and missing value imputation:
\begin{lstlisting}
df_filtered = df_trans[(df_trans['status'] == 'SUCCESS') & (df_trans['amount'] >= 5_000_000)].copy()
df_filtered['fee'] = df_filtered['fee'].fillna(0)
\end{lstlisting}
  \textit{Rubric:} Correct boolean filtering with bitwise \texttt{\&}: 0.25 pts; Correct \texttt{.fillna(0)}: 0.25 pts.

  \item \textbf{(1.0 pts):} Groupby and aggregation expression:
\begin{lstlisting}
df_branch_summary = df_filtered.groupby('branch').agg(
    total_amount=('amount', 'sum'),
    avg_fee=('fee', 'mean'),
    trans_count=('amount', 'count')  # Or 'size'
).sort_values(by='total_amount', ascending=False)
\end{lstlisting}
  \textit{Rubric:} Correct \texttt{groupby('branch')}: 0.25 pts; Correct aggregations (\texttt{'sum'}, \texttt{'mean'}, \texttt{'count'}): 0.5 pts; Correct \texttt{sort\_values(..., ascending=False)}: 0.25 pts.
\end{enumerate}

\subsection*{Question 3: Multi-plot Layouts with Matplotlib (1.0 pt)}
\begin{enumerate}[label=\textbf{\alph*)}]
  \item \textbf{(0.5 pts):} Statement to initialize subplots:
\begin{lstlisting}
fig, axes = plt.subplots(1, 2, figsize=(12, 4))
\end{lstlisting}
  \textit{Rubric:} Correct grid specification \texttt{(1, 2)}: 0.25 pts; Correct \texttt{figsize=(12, 4)}: 0.25 pts.

  \item \textbf{(0.5 pts):} Plotting and layout commands:
\begin{lstlisting}
axes[0].plot(df_kpi['month'], df_kpi['profit'], color='blue', marker='o')
axes[1].bar(df_kpi['month'], df_kpi['profit'], color='steelblue')
plt.tight_layout()
\end{lstlisting}
  \textit{Rubric:} Correct line plot on \texttt{axes[0]} with circular marker: 0.2 pts; Correct bar chart on \texttt{axes[1]}: 0.2 pts; Correct \texttt{plt.tight\_layout()}: 0.1 pts.
\end{enumerate}

\subsection*{Question 4: Exploratory Visualization \& Business Insights (EDA) (1.5 pts)}
\begin{enumerate}[label=\textbf{\alph*)}]
  \item \textbf{(0.75 pts):}
  \begin{itemize}
    \item \textbf{Chart type:} **Box Plot** (or Violin Plot). \textit{(0.25 pts)}
    \item \textbf{Seaborn statement:} \textit{(0.5 pts)}
\begin{lstlisting}
sns.boxplot(data=df_cards, x='is_fraud', y='tx_amount', hue='card_type')
# (Acceptable variant: x='card_type', y='tx_amount', hue='is_fraud')
\end{lstlisting}
  \end{itemize}

  \item \textbf{(0.75 pts) --- Analytical Takeaway \& Policy Recommendation:}
  \begin{itemize}
    \item \textit{Analytical Takeaway (0.375 pts):} Credit card fraud typically targets high-credit limits with anomalous transaction sizes (extreme outliers between 50M and 200M VND), contrasting sharply with regular transactions ($\le 10$M VND) and debit card behavior.
    \item \textit{Policy Recommendation (0.375 pts):} The bank should establish automated step-up authentication (Smart OTP/biometrics) or immediate transaction freezes for all single credit card charges exceeding 50 million VND.
  \end{itemize}
\end{enumerate}

\begin{center}
\textbf{------------------- END OF SOLUTIONS -------------------}
\end{center}

\end{document}
